%to have line numbers
%\RequirePackage{lineno}
\documentclass[10pt, letterpaper]{article}      
\usepackage[margin=.1cm,font=small,labelfont=bf]{caption}[2007/03/09]
%\usepackage{endnotes}
%\let\footnote=\endnote
\usepackage{setspace}
\usepackage{longtable}                        
\usepackage{anysize}                          
\usepackage{natbib}                           
%\bibpunct{(}{)}{,}{a}{,}{,}                   
\bibpunct{(}{)}{,}{a}{}{,}                   
\usepackage{amsmath}
\usepackage[% draft,
pdftex]{graphicx} %draft is a way to exclude figures                
\usepackage{epstopdf}
\usepackage{hyperref}                             % For creating hyperlinks in cross references
\hypersetup{pdfborder={0 0 0.4}} %have nice light boxes on refs

% \usepackage[margins]{trackchanges}

% \note[editor]{The note}
% \annote[editor]{Text to annotate}{The note}
%    \add[editor]{Text to add}
% \remove[editor]{Text to remove}
% \change[editor]{Text to remove}{Text to add}

%TODO make it more standard before submission: \marginsize{2cm}{2cm}{1cm}{1cm}
\marginsize{1cm}{1cm}{.5cm}{.5cm}%{left}{right}{top}{bottom}   
					          % Helps LaTeX put figures where YOU want
 \renewcommand{\topfraction}{1}	                  % 90% of page top can be a float
 \renewcommand{\bottomfraction}{1}	          % 90% of page bottom can be a float
 \renewcommand{\textfraction}{0.0}	          % only 10% of page must to be text

 \usepackage{float}                               %latex will not complain to include float after float

\usepackage[table]{xcolor}                        %for table shading
\definecolor{gray90}{gray}{0.90}
\definecolor{orange}{RGB}{255,128,0}

\renewcommand\arraystretch{.9}                    %for spacing of arrays like tabular

%-------------------- my commands -----------------------------------------
\newenvironment{ig}[1]{
\begin{center}
 %\includegraphics[height=5.0in]{#1} 
 \includegraphics[height=3.3in]{#1} 
\end{center}}

 \newcommand{\cc}[1]{
\hspace{-.13in}$\bullet$\marginpar{\begin{spacing}{.6}\begin{footnotesize}\color{blue}{#1}\end{footnotesize}\end{spacing}}
\hspace{-.13in} }

%-------------------- END my commands -----------------------------------------



%-------------------- extra options -----------------------------------------

%%%%%%%%%%%%%
% footnotes %
%%%%%%%%%%%%%

%\long\def\symbolfootnote[#1]#2{\begingroup% %these can be used to make footnote  nonnumeric asterick, dagger etc
%\def\thefootnote{\fnsymbol{footnote}}\footnote[#1]{#2}\endgroup}	%see: http://help-csli.stanford.edu/tex/latex-footnotes.shtml

%%%%%%%%%%%
% spacing %
%%%%%%%%%%%

% \abovecaptionskip: space above caption
% \belowcaptionskip: space below caption
%\oddsidemargin 0cm
%\evensidemargin 0cm

%%%%%%%%%
% style %
%%%%%%%%%

%\pagestyle{myheadings}         % Option to put page headers
                               % Needed \documentclass[a4paper,twoside]{article}
%\markboth{{\small\it Politics and Life Satisfaction }}
%{{\small\it Adam Okulicz-Kozaryn} }

%\headsep 1.5cm
% \pagestyle{empty}			% no page numbers
% \parindent  15.mm			% indent paragraph by this much
% \parskip     2.mm			% space between paragraphs
% \mathindent 20.mm			% indent math equations by this much

%%%%%%%%%%%%%%%%%%
% extra packages %
%%%%%%%%%%%%%%%%%%

\usepackage{datetime}


\usepackage[latin1]{inputenc}
\usepackage{tikz}
\usetikzlibrary{shapes,arrows,backgrounds}


%\usepackage{color}					% For creating coloured text and background
%\usepackage{float}
\usepackage{subfig}                                     % for combined figures

\renewcommand{\ss}[1]{{\colorbox{blue}{\bf \color{white}{#1}}}}
\newcommand{\ee}[1]{\endnote{\vspace{-.10in}\begin{spacing}{1.0}{\normalsize #1}\end{spacing}\vspace{.20in}}}
\newcommand{\emd}[1]{\ExecuteMetaData[/tmp/tex]{#1}} % grab numbers  from stata

%TODO before submitting comment this out to get 'regular fornt'
\usepackage{sectsty}
\allsectionsfont{\normalfont\sffamily}
\usepackage{sectsty}
\allsectionsfont{\normalfont\sffamily}
\renewcommand\familydefault{\sfdefault}

%\usepackage[margins]{trackchanges}
\usepackage{rotating}
\usepackage{catchfilebetweentags}

\usepackage{abstract}
\renewcommand{\abstractname}{}    % clear the title
\renewcommand{\absnamepos}{empty} % originally center
%-------------------- END extra options -----------------------------------------
\date{{}\today \hspace{.2in}\xxivtime}
\title{  % remember to have Vistula University!!
  Geographies of Wellbeing in Poland\\ (Urban-Rural, Regions, and Cities) 
}
\author{
% Adam Okulicz-Kozaryn\thanks{EMAIL: adam.okulicz.kozaryn@gmail.com
%   \hfill I thank XXX.  All mistakes are mine.} \\
% {\small Rutgers - Camden  % and Vistula University
% }
}

\begin{document}

%%\setpagewiselinenumbers
%\modulolinenumbers[1]
%\linenumbers

\bibliographystyle{/home/aok/papers/root/tex/ecta}
\maketitle
\vspace{-.4in}
\begin{center}

\end{center}


\begin{abstract}
\noindent
\end{abstract}
\vspace{.15in} 
\noindent{\sc %XXX TODO add to ebib as keyword PAPER-CODE-NAME and tag with ebib keywords 
}
\vspace{.25in} 

\vspace{-2.5in} 


\begin{spacing}{1.4} %TODO MAYBE before submission can make it like 2.0
\rowcolors{1}{white}{gray90}

\section{Wellbeing}

The proposed study's outcome of interest is multidimensional human
wellbeing or flourishing. Two key dimensions are happiness or life
satisfaction\footnote{We use the terms interchangeably as is the usual practice
  in the social indicators literature, but specifically we mean global cognitive
 evaluation of life as whole rather than momentary and fleeing affective
 happiness.} and freedom (agency/autonomy).\footnote{Specifically we mean
 so called \texttt{freedom to}, not \texttt{freedom from} \citep{aok_free_from_to}.
Poles are \texttt{free from}: live in a free country, have no coercion,
 and are free from restrictions and impediments. But are they \texttt{free to}? 
 Are Poles actually able to choose, control and direct one's own life
 \citep{ee_ls15}?
% Not so much \citep{ee_ls15}.
}

Happiness is now recognized as a development goal in
addition to GDP--the UN has passed resolution A/RES/65/309: ``Happiness: towards a holistic approach to
development'' stating that ``the pursuit of happiness is a fundamental human
goal'' and ``recognizing  that  the  gross  domestic  product  indicator [GDP]
by  nature  was  not  designed to and does not adequately reflect the happiness
and well-being of people.''\footnote{See
  \url{https://docs.un.org/en/A/RES/65/309}. Also see report promoting
  wellbeing as a human development measure by Nobel prize laureates Sen and Stiglitz \citep{stiglitz09al}.}

Freedom (agency/autonomy) is a human right as stated by the  Universal
Declaration of Human Rights (art. 26, 29): ``freedom for the realization of one's human potential.'' 
Freedom  is an end-in-itself%\citep{aok_free_from_to}
, as recognized by 
\citet{fromm13,fromm12,fromm92,fromm64,fromm44,fromm94} and \citet{sen00}.

% \textbf{do check if this needs to be blind, then need to rewrite}
% I have coined the term ``urban-rural happiness gradient'' \citep{aok11a}
% % and the new line of literature has started.
% I am the leader of this field.\footnote{For instance see the discussion in framing essay IV in
%   \citet{johansen22}.}
There is a new geography of wellbeing field \citep[e.g.,][]{ballas2013geography}. 
 The field has mostly international focus across
countries, with only few recent studies focusing on specific
countries.\footnote{See next section.} 
 %

Now the time is right for a geography of wellbeing study in Poland.
In 2026 Poland is not very different
from Western Europe anymore. Poland is a quite developed country now. Indeed, Polish PCGDP
PPP\footnote{Per Capita Gross Domestic Product Purchasing Power Parity adjusted.} 
is now about the same as in Japan, and Polish minimum wage of 31 PLN or 8.15 USD
(as of 2026) is higher than the US minimum wage at 7.25 USD.  
%https://en.wikipedia.org/wiki/List_of_countries_by_GDP_(PPP)_per_capita

 It was not like that before, however.  
 I remember about 10 years ago I gave a talk on happiness in Warsaw, 
%vistula, maybe little more recently like 2019
 and students didn't seem to see a need to care about happiness. Traditional
topics like economic growth and personal income seemed more important. 

Wellbeing is increasingly relevant as Poland has been transitioning from
 developing country
  into a developed country so quality of life and wellbeing
increasingly matter over economic/survival basics such as income and
unemployment in the past. 
 When Poland was quite poor the focus rightly so was on survival
  and economic growth.  As Poland has
experienced economic miracle over past several years,\footnote{For instance Warszawski
Index Gieldowy (WIG) has increased about 120\% over past 5 years.} now it's
time to focus on wellbeing. 
 %
For visualization of the changing values from survival to wellbeing along with
the economic growth see figure \ref{ingle}.

\begin{figure}[H]
 \includegraphics[height=2.3in]{/home/aok/papers/root/old/2025/swbColCitReg/tex/ingle.pdf}\centering
\caption{\label{ingle} From survival to wellbeing as economic growth increases \citep{inglehart97}.}
\end{figure}


\section{Geographies: Urban-Rural, Regions, and Cities}

%!!!!!!!!!!!!!!!!!!!!!!!!!!!!!!!!!!!!!
%$>>>$guess cp here lit from my papers esp recent ones:\\
%(make sure i caover everything i have in hypotheses!!!)
%guess moree qol livability stuff and shrinking gorwing\\



%https://en.wikipedia.org/wiki/List_of_European_countries_by_area
%https://en.wikipedia.org/wiki/List_of_European_countries_by_population
 Poland is a medium-large country, 9th both in terms of area and population in Europe,\footnote{
 See \url{unstats.un.org} and \url{https://population.un.org}.}
 and yet we know little about geographies of wellbeing in Poland and lack a
dedicated comprehensive study. Hence, I propose this study to fill the gap.   
 %
Poland is a diverse country in terms of its urban-rural, regional, and city geographies. 

%!!!!!!! MAYBE add there that diagram from Ott chapter
Urban-rural is not a binary, but a gradient--we will use several steps to
 capture this gradient. Urban-rural gradient matters a great deal, more than
most scholars realize. Indeed, it is difficult to overestimate its importance--the key
point is that urbanicity is not only a built environment, but a way of life \citep{wirth38}.%--details in the data section.
 %

%Several urban-rural happiness gradient studies are worth mentioning. 
 \citet{weckroth21} is a closely related study to our proposed study--they analyze SWB by urban-rural
using ESS in Europe--we will also use ESS among other datasets, 
 but focus on Poland and add more recent waves and more wellbeing indicators. 
 %
\citet{knies26} is another close ESS study\footnote{And there are 2 more ESS
  studies: \citet{ivanovic2022subjective} is an urban-rural study, but again
only 2 waves and no focus on Poland. \citet{vaznoniene2019wellbeing} is
interesting ESS study in a sense that it does focus on one country and discusses
country specific intricacies like us, but it is Lithuania. \citet{aok21} is a major recent study using WVS, but as most research it focuses on cross-country
 patterns, not Poland.} --it uses multiple indicators, but only one wave 6 and like other
research does not focus on Poland. %We will use all 12 waves.
 %
There are no studies on urban-rural freedom gradient, not just in Poland but
anywhere--ours is the first.\footnote{\citet{yamagishi12} studies urban-rural
  freedom but in very different way from us--they show that urbanites are more
  likely to chose extravagant pens than rural folks. We will study freedom for
  urban-rural and regions (we do not have 
freedom variable for cities).} Next we turn to regions.
 
 Regions are an useful level of analysis as there are
 great differences within countries and across regions, and Poland is no
exception.  There are life satisfaction studies across regions, but no study focuses on Poland \citep[e.g.,][]{aok10_geo,ballas23,ziogas23,ziogas20}.
Similarly, quality of life and wellbeing studies across regions and cities focus
on cross-country patterns, not Poland, and tend to use older datasets \citep{aok-swbLivability18,aok13liavbility}.
Poland has 16 voivodeships, so called NUTS2 level
regions.\footnote{Voivodeships$=$NUTS2, except  Mazowieckie (the Warsaw
  region)--1 single voivodeship, but 2 NUTS2: Warsaw Metropolitan and Mazovian
  Region; so total of 17 NUTS2 regions.} %(in European statistical terms).
 At NUTS2 level we will investigate geographic clustering and the relationship
 between objective development measures such as income and subjective measures (SWB). 
 %
There are also interesting population dynamics at NUTS2 level, for
instance,  Mazovia (where Warsaw is located) has been growing, but most  
other regions have been stable or shrinking. We will investigate the effect of
population change on wellbeing--we know that population dynamics are highly
correlated with SWB \citep{growShrink25}, but again, we do not know the
relationship in Poland. 
%https://en.wikipedia.org/wiki/List_of_cities_and_towns_in_Poland

Finally, we have data representative of 
Bialystok, Gdansk, Cracow, and Warsaw with domain satisfactions (life, place,
personal, and financial), and various quality of life indicators regarding
amenities (e.g., public goods and services such as parks)  and
disamenities (e.g., community problems such as crime). We will investigate 
these across the cities.    
 %
 % For instance, Warsaw has been clearly growing, while other cities such as
 % Poznan have been mostly flat, and yet others such as Lodz and Szczecin have
 % been shrinking.
 %
 Specifically, large Warsaw (and perhaps Cracow) is likely to differ from small
 Bialystok. 
 Warsaw stands out as not only the largest city, but also possibly a primate
 city, disproportionately larger and more influential than any other city.\footnote{At
$>3m$ Warsaw metro is larger than Katowice metro at about 2.5m and third
Karkow metro at 1.5m.
 %
While Katowice metro is not half or less of Warsaw metro as
 ``the law of the primate city''
 would indicate \citep[see discussion in][]{aok12zipf}, it is only due to the
 location of Katowice in regional industrial cluster (Gornoslaski Okreg
 Przemyslowy (GOP)), Katowice city itself is much smaller than  Warsaw at
 0.28m v 1.9m. 
 %
Maybe Warsaw is not such clear primate as Cairo  \citep{ebshoyEgypt24}, but it
could be similar in its primate status to Bogota \citep{unhappyBog26}.
}
 
Last but not least, there are several related Polish studies in smaller local
journals that are worth mentioning as they are mostly about Poland
specifically--again, the international literature tends not to focus on Poland.
 %The closest first.
 \citet{sompolska2020quality} is quite close, but uses different data: Polish
Social Diagnosis dataset and does not control for SWB predictors. 
 %
\citet{vaznoniene2021territorial} uses ESS and studies urban-rural SWB like us, but
 focuses on Lithuania v Poland, and uses waves 4 and 9 only and 3 urbanicity
 categories only, and does not control for SWB predictors. \citet{wojewodzka2023rural}
 is the same, just uses waves 2 and 8. 
  %
\citet{tobiasz2017urban} does the analyses across the regions and urban-rural,
 %and while social connection and loneliness are closely related to SWB,
 but their sample is elderly and the study is about health only.
 \citet{stangierska2022urban} studies quality of life across green v not green
 urban areas, but only across Warsaw neighborhoods, and the sample size is small at 381 persons.
 \citet{sompolska2021happier} uses Polish Social Diagnosis and studies urban-rural 
 SWB--but separately: predictors of SWB in urban separately from predictors of
 SWB in rural. %--it is different from our study.
 Overall, while there are studies focusing on Poland, like our proposed study, they have
 important limitations as summarized above.\footnote{ 
 And there are several studies that differ from the
 above. \citet{churski2022geographical} has very fine 
 geographical resolution of 380 units across Poland, but they do not study SWB,
  only the related objective indicators such as jobs. 
 \citet{weziak2025understanding} is one of the very many of the studies
 co-authored with VanderWeele, this time in Poland, but there is no geography in
 the study.  
 %
 \citet{wkeziak16} uses the same flash barometer 366 that we will use, but does
 not focus on Poland and instead on all European cities. 
}
 Hence,  we know little about wellbeing geographies in Poland--we
 turn to our, yet un-answered, hypotheses.  

\section{Hypotheses}

% The research question: 
% This research is about geography of wellbeing in Poland (urban-rural gradient, regions, and cities).

% The hypotheses listed below are a primary set, we will be refining them and
% adding other hypotheses at the literature review stage.

\noindent %Hypotheses
%(SWB$=$(multidimensional) Subjective Wellbeing)):\\
$H_1$ Cities have lower multidimensional Subjective Wellbeing (SWB) than rural areas, with the largest
drop occurring for the largest places. [WVS, ESS, FB539]\footnote{WVS$=$World Values
  Survey; ESS$=$European Social Survey; FB539$=$Flash Barometer \#539--these and other datasets are described
  in the next section.}\\ %ess
$H_1a$ Certain social groups such as working-age adults and higher social classes
have higher SWB in cities v rural areas. [WVS, ESS]\\
$H_1b$ Over time, from 1989 to 2012, the SWB gap between cities and rural
places has increased. [WVS]\\ 
$H_2$ More remote, less populated, and lower density regions have higher SWB than
more central, more populated and higher density regions. [ESS, FB539]\\ %flash b: just maps no regressions too few units
$H_2a$ Regions that experienced population growth will have higher SWB than flat
or shrinking regions. [ESS, FB539]\\ 
$H_2b$ Across regions correlation between objective quality of life measures such as income with
SWB measures will be moderate only at $<0.5$  thus revealing inadequacy of
  traditional progress/development measures such as
income. [ESS, FB539]\\
$H_2c$ Urban-rural happiness gradient across regions will be relatively stable,
i.e., larger places have higher SWB than smaller places. [ESS, FB539]\\ 
$H_3$ The relationship between quality of life (local amenities and disamenities) and domain satisfactions (life, place, personal, and financial) will differ across                                  
 Bialystok, Gdansk, Cracow, and Warsaw. [FB366,419] %\citep{aok-swbLivability18}

 
 
\section{Data and Measurement}
%maybe that diagnoza spoleczna--or may turn out its crap, maybe someone already said so

%\subsection{Multidimensional Wellbeing}
The goal is to do a comprehensive study of Polish geographies of wellbeing, hence
we use multiple datasets, each having its own strengths: World Values Survey (WVS), European Social Survey
(ESS), and Flash Barometers (FB) for regions (FB539) and cities (FB366,419).

WVS has a usual SWB item:
 ``All things considered, how satisfied are you with your life as a whole these days?''
 on scale 1-10. 
Such SWB measures tend to be reasonably valid and reliable
%  \citep{myers00,ditella06m,layard05,bray06g,sandvik93ds,clark08al,diener09}.
\citep[][ch. 5]{diener09}. 
% 
WVS allows us to examine the geographies of wellbeing over time, which is important
  as Poland has evolved greatly, and WVS data goes back to 1989. We have 4k observations, about
  1k per wave.\footnote{We use waves 1989, 1997, 2005, and 2012.
 For the 1997, 2005, and 2012 waves we have 8
rural-urban categories -2k, 2-5k, 5-10k, 10-20k, 20-50k, 50-100k, 100-500k, and
500k-. In 1989 we only have information about 49 voivodeships, but
 urbanicity can be inferred from this information indirectly,
 and proximity to a major city directly.} 
 Also, an important advantage of WVS is
 that urbanicity is recorded by an interviewer.\footnote{It is more objective than subjective respondent self-reports. For
   instance, respondents may not know well whether they live in a 20k or 30k place.}

Freedom is measured in WVS,\footnote{"Some people feel they have completely free choice and control over their lives, while other
      people feel that what they do has no real effect on what happens to them. Please use this scale
      where 1 means 'none at all' and 10 means 'a great deal' to indicate how much freedom of choice
      and control you feel you have over the way your life turns out." 
    Validity and reliability of the measure seem fine and the variable has been used successfully
 in the literature \citep[e.g.,][]{verme09,inglehart08,veenhoven00_b,aok_free_from_to}. }
%if anything it could be interpreted slightly differently in different
%languages/cultures but that is not a problem here for Poland only.}
 and in some ESS waves 
 %which has several waves %3,6,12
 containing a whole battery of wellbeing items measuring functionings and vitality. 
 %
Preliminary analyses using ESS 1-6 cumulative file already had over 10k
observations for Poland.\footnote{The data will be updated upto wave 12 that will be
available in 2027.}
  %
 We will use ESS items such as those in wave 6 on functioning ``free to decide how to live my life,''
  %``enthusiastic about what you are doing,''
 and vitality ``had lot of energy.''\footnote{ESS urbanicity is measured with: a big city, suburbs or outskirts
of big city, town or small city, country village, farm or home in
countryside. Per wellbeing see overview at 
  \url{https://www.europeansocialsurvey.org/findings/europeans-wellbeing/measuring-wellbeing}. Wave 
6 wellbeing items:
\url{https://europeansocialsurvey.org/sites/default/files/2026-01/ESS12-wellbeing-module-design-template.pdf}
 see Table 1: Items from the wellbeing module grouped by the dimension of
 wellbeing they relate to. We also will use new wave 12 that will be available in 2027 and will
 contain a wellbeing module. See \url{https://europeansocialsurvey.org/sites/default/files/2026-01/ESS12-wellbeing-module-design-template.pdf}.}
 

FB539 (2024) has 5.2k observations and is representative of
 17 NUTS2 regions, and has a usual SWB
 question similar to WVS.\footnote{We may also use
  the earlier flash barometers: \#497 (2021), \#472 (2018), \#427 (2015),  \#356 (2012)
  depending on how strong and robust the latest wave results are.} 
% \begin{figure}[H]
%  \includegraphics[height=2in]{nuts2.png}\centering
% \caption{\label{nuts2} 17 regions (voivodeships) so called NUTS2 regions}
% \end{figure}
 %
FB366,419 (2012, 2015) are representative of 4 Polish cities:
 Bialystok, Gdansk, Cracow, and Warsaw. The sample size is about 2000
 observations per wave, and again we have a usual SWB question like in WVS.
%later can also try Flash Eurobarometer 156 (Urban Audit Perception Survey 2004)

 
\section{Models}


We will mostly use ordinary least squares (OLS).\footnote{Although OLS assumes cardinality of the
outcome variable, and SWB measures are technically ordinal, OLS is an appropriate estimation method.
\citet{carbonell04} has shown that OLS results are substantially the same as
those from discrete models, and  OLS has become the default method in happiness
research \citep{blanchflower11}. Theoretically, while there is still debate
about the cardinality of SWB, there are strong arguments to treat it as a
cardinal variable \citep{ng96,ng97,ng11}.} For robustness/sensitivity
checks for the key models we will use multinomial logit (as the ordinal logit
tends to fail proportional odds assumption). As another robustness/sensitivity
check we are considering machine learning models such as gradient boosting
models \citep{oparina25}, specifically XGBoost
 (\url{https://xgboost.readthedocs.io}). 

%\section{controls}
We will use a usual set of controls such as those in the literature: 
for life satisfaction \citep[e.g.,][]{aokcities,aokrel}; and for freedom
\citep[e.g.,][]{pitlik2016free,nikolaev2016give}. 
 %
 A key robustness/sensitivity check will be specification curve analysis (SCA)
\citep{simonsohn20}. And %, related,
 we will pay extra attention to the choice 
 of control variables and avoid overcontrol bias.\footnote{
%\begin{quote}
``What we really need, to identify/estimate X$->$Y, is to control for other
determinants of Y that are also antecedents of X (so, W$->$X). If we include
controls where the relationship goes the other way (X$->$W), we will exacerbate
bias in our estimate of X$->$Y. (The relevant term for that situation is
'overcontrol bias')'' \citep[][p. 5]{bartram24}.
%\end{quote}
}

To examine geographic patterns among regions we will use choropleth maps,
Moran's I, and Local Indicators of Spatial Association (LISA) \citep{anselin1995local}.\footnote{ 
We may do some multilevel modeling (MLM) at NUTS2 level as in
\citet{weckroth21} but focusing on Poland. Yet MLM is problematic with only 17
U/As at Level 2--if we use it, results will be postponed to appendix. 
We may use Small Area Estimation \citep{moretti24}. 
}



\section{Significance, Policy, and Practice}

%guess third of points is: Potencjalne znaczenie Projektu dla rozwoju nauki

The proposed research will contribute to the literature and produce new
knowledge.   
 We will write several papers in this project, and through further research
mentoring and leadership additional ones.
%See for details in the application form under ``Dzialania mentoringowe.'' 
%\textbf{see if thats where it fits}

But the research also has policy and practical implications. First, as stated in the
introduction, the time has come for Poles to focus on wellbeing as opposed to
just survival. We have modernized infrastructure, higher incomes, and booming economy. The %past 
focus on survival and economy
 can now give a way to 
 %now shift towards
 wellbeing and flourishing.
  With the proposed study we aim to aid in this pursuit of happiness. 
 This is called evidence-based pursuit of happiness in the literature \citep{veenhoven12}.\footnote{There is a
  psychological nuance here--one should not focus fully on the pursuit of
  happiness and chase it or value it intensely
  \citep{mauss11,gruber11}, neither be concerned about it and judge it
  \citep{zerwas24}--it may backfire. Rather, one needs to know what is important
for wellbeing from the academic literature, like the proposed study, and include
this information in everyday decision-making. For instance, humans should
 buy experience such as travel and time with others, not things such
as conspicuous consumption items \citep{aok_ls_car15,veblen05a,veblen05b,thoreau95}.}
 The proposed study will find out which parts of Poland enjoy high wellbeing and for
 whom--which social groups, e.g., by age and social class, flourish where? 
 Our findings should help Poles increase their wellbeing.

 SWB can and should be used for public policy and administration \citep{diener09}.
% provides a good discussion of why potential problems with happiness are not serious
% enough to make it unusable for public policy--see especially  ch. 6.  
% These potential problems include  genetic
% determination of happiness \citep{lykken96t},  adaptation
% \citep{brickman78cj},  and various comparisons \citep{michalos85}.}  
% 
 SWB can serve as a ``yardstick'' \citep{veenhoven88}--arguably the ultimate outcome of public
 policy and administration is SWB.\footnote{While the idea of using wellbeing data for policy and
    administration is new in Poland, it is not new in general
    \citep{stiglitz09al,diener09}.} "If you can't measure it, you can't improve
 it"--if governments don't measure wellbeing, they can't improve wellbeing.  
  I will reach out to Polish government to implement the study findings in
  practice.\footnote{One simple implementation is simply to educate policy
    makes, public administrators, and citizenry about the wellbeing metrics and
    study's findings so that we can make more informed decisions. Another
    example is to use the wellbeing metrics to evaluate policy. For instance, if
  there a major change in policy or investment in public goods, we can actually
  measure whether it produced an increase in wellbeing by using before-after
  comparative group research design (difference-in-difference).}
  I plan to start locally with Wielkopolska voivodeship,
  where I will be physically carrying out the study, and other regional
  governments may be involved as well. %as I will be carrying out the study.
  Likewise, some city
  governments may be involved as the study will use data representative of 4
  cities.
  
  




  

  

%  instead \ExecuteMetaData[../out/tex]{ginipov} do \emd{ginipov}

% \begin{figure}[H]
%  \includegraphics[height=3in]{../out/gov_res_trust.pdf}\centering\label{gov_res_trust}
% \caption{woo}
% \end{figure}


%TODO !!!! have input here aok_var_des



% %table centered on decimal points:)
% \begin{table}[H]\centering\footnotesize
% \caption{\label{freq_im_god} importance of God}
% \begin{tabular} {@{} lrrrr @{}}   \hline 
% Item& Number & Per cent   \\ \hline
% 1(not at all)&    9,285&  9\\
% 2&    3,555&        3\\
% 3&    3,937&        4\\
% 4&    2,888&        3\\
% 5&    7,519&        7\\
% 6&    5,175&        5\\
% 7&    6,050&        6\\
% 8&    8,067&        8\\
% 9&    8,463&        8\\
% 10&   52,385&       49\\
% Total&  107,324&      100\\ \hline
% \end{tabular}\end{table}


% % Define block styles
% \tikzstyle{block} = [rectangle, draw, fill=black!20, 
%     text width=10em, text centered, rounded corners, minimum height=4em]
% \tikzstyle{b} = [rectangle, draw,  
%     text width=6em, text centered, rounded corners, minimum height=4em]
% \tikzstyle{line} = [draw, -latex']
% \tikzstyle{cloud} = [draw, ellipse,fill=black!20, node distance = 5cm,
%     minimum height=2em]
    
% \begin{tikzpicture}[node distance = 2cm, auto]
%     % Place nodes
%     \node [block] (lib) {liberalism, egalitarianism, welfare};
%     \node [block, below of=lib] (con) {conservatism, competition, individualism};
%     \node [cloud, right of=con] (ls) {well-being};
%     \node [block, below of=ls] (cul) {genes, culture};
%     \node [b, left of =lib, node distance = 4cm] (c) {country-level};
%     \node [b, left of =con,  node distance = 4cm] (c) {person-level};
%     % Draw edges
%     \path [line] (lib) -- (ls);
%     \path [line] (con) -- (ls);
%     \path [line,dashed] (cul) -- (ls);
% \end{tikzpicture}


%PUT THIS NOTE, polish and put to /root/author_what_data --ALWAYS
%stick here stuff as i run it!!! maybe comment out later...


%!!!!!!!!!TODO: see cityAgency for cool setup :)))

% \section*{\Huge ONLINE APPENDIX}
% \textbf{[note: this section will NOT be a part of the final version of
%   the manuscript, but will be available online instead]} %hence everything below
%                                 %is organized byu section, not subsection
% !!!
% have most of the stuff outputted to online appendix:)--start with that and then
% select stuff to paper--have brief narrative describng patterns in online app too
% !!!

% \section*{Variables' definitions, coding, and distributions}
% \label{app_var_des}


% %\input{/tmp/a.tex} %aok_var_des

% % \begin{spacing}{.9}
% %   \begin{table}[H]\centering \caption{Summary statistics.} \label{sumSta} \begin{scriptsize} \begin{tabular}{p{1.8in}p{.5in}p{.5in}p{.5in}p{.5in}p{.5in}p{.5in}p{.5in}p{.5in}p{.5in}p{.5
% %             in}p{.5in}p{.5 in}}\hline
% %         \input{/tmp/aha2.tex}
% %          \end{tabular}\end{scriptsize}\end{table}
% % \end{spacing}

% % \begin{spacing}{.9}
% %   \begin{table}[H]\centering \caption{Correlation matrix.} \label{sumSta} \begin{scriptsize} \begin{tabular}{@{}
% %           p{1.2in} rrrrrrrrrrrrr @{}}\hline
% %         \input{/tmp/ahb2.tex}\hline
% %          \end{tabular}\end{scriptsize}\end{table}
% % \end{spacing}



% Table XXX shows variable distributions. If a variable has more than
% 10 categories it is classified into bins...

% %\input .... %TODO !!!! have input here histograms

% \section*{Additional Descriptive Statistics}
% \label{app_des_sta}

% %make sure i have [H] or h! ???
% % \begin{table}[H]
% % \caption{}
% % \centering
% % \label{}
% % \begin{scriptsize}
% % \input{../out/reg_c.tex}
% % \end{scriptsize}
% % \end{table}

%\newpage
%\theendnotes


% \begin{scriptsize}
% \begin{spacing}{.9} 
% \begin{verbatim}
% . tabstat ls, by(lwd7) stat(mean sd) format(%9.1f)

% Summary for variables: ls
% Group variable: lwd7 (RECODE of lwd (do what you like))

%   lwd7 |      Mean        SD
% -------+--------------------
%    0-4 |       5.6       2.7
%      2 |       7.6       2.3
%      3 |       7.6       1.8
%      4 |       8.0       1.5
%      5 |       7.8       1.9
%      6 |       8.2       1.5
%      7 |       8.9       1.5
% -------+--------------------
%  Total |       8.3       1.9
% ----------------------------
% \end{verbatim}
% \end{spacing}{.9}
% \end{scriptsize}


\bibliography{/home/aok/papers/root/tex/ebib.bib,powroty.bib,/home/aok/papers/cityAgency/tex/cityAgency.bib}

\end{spacing}
\end{document}
